\section{Wnioski}
\label{sec:conclusions}

Praca odpowiada na 4 pytania z wprowadzenia.

\paragraph{Pytanie (i): model.}
Język jest parą: zbiór zdań i relacja interpretacji. Reguła PK jest ograniczeniem albo
konwencją; obie operacje są przecięciem z ustalonym zbiorem, więc reguły są przemienne.
PK dąży do tego, żeby relacja interpretacji była funkcją. Numerowane reguły PK 0.1 tego
celu nie gwarantują: praca wskazuje 2 zdania zgodne z regułami, które mają 2 odczytania.

\paragraph{Pytanie (ii): koszty.}
Teza z tytułu zachodzi w 3 znaczeniach i nie zachodzi w 1 znaczeniu.
\begin{itemize}[leftmargin=1.6em,itemsep=1pt]
  \item PK zmniejsza zbiór odczytań. Zdanie z 8 frazami przyimkowymi ma \pkAttachEight{}
  odczytania; zdanie PK ma 1 odczytanie. Brak informacji w zdaniu PK wynosi 0 bitów.
  \item PK zmniejsza koszt parsowania w modelu gramatyki: zmierzony wykładnik spada
  z \pkParseExpNat{} do \pkParseExpPk.
  \item PK wyznacza granicę kosztu wnioskowania. Zbiór wymagań bez przeczeń ma kontrolę
  spójności w czasie liniowym. Zbiór wymagań z 1 przeczeniem w zdaniu ma kontrolę spójności
  NP-zupełną.
  \item PK nie zmniejsza kosztu przejścia tekstu przez model językowy. Tekst PK jest
  dłuższy: mediana narzutu wynosi \pkTokMedian{} \% tokenów.
\end{itemize}

\paragraph{Pytanie (iii): pomiar.}
Modele toru Claude odczytują zdania PK zgodnie z kluczem w co najmniej
\pkAccPkClaudeMin{} \% przypadków. Modele lokalne osiągają od \pkAccPkLocalMin{} \%
do \pkAccPkLocalMax{} \% włącznie przy poziomie przypadku 50 \%. Zdanie zwykłe dzieli
modele po kontroli w \pkNatSplit{} \% pozycji. Korzyść z PK zależy więc od czytelnika:
PK pomaga modelowi, który czyta polskie zdanie poprawnie, i nie zastępuje tej zdolności.
Próg opłacalności w modelu kosztu wynosi \pkBreakEvenStar{} \% błędnych odczytań zdania
zwykłego.

\paragraph{Pytanie (iv): hipoteza Sapira-Whorfa.}
PK zmienia listę kategorii, które autor wyraża obowiązkowo: wykonawca, ilość, granica,
chwila, źródło wiedzy. Dla człowieka praca stawia hipotezę: pisanie w PK ujawnia autorowi braki
w wiedzy o opisywanej czynności. Dla modelu językowego praca rozdziela hipotezę na 3 tezy.
Zależność odpowiedzi od tekstu zachodzi z budowy modelu. Efekt rejestru zachodzi dla
modeli toru Claude. Efekt nawyku widać w pozycjach typu konwencja: konwencja niezgodna
z nawykiem działa dopiero z legendą.

\paragraph{Zalecenia dla standardu.}
Wyniki dają 4 zalecenia dla kolejnej wersji PK.
\begin{enumerate}[leftmargin=1.8em,itemsep=1pt]
  \item Dodać regułę o zasięgu zakazu nad spójnikiem \pkq{i} (luka G1).
  \item Dodać regułę o zasięgu 2 kwantyfikatorów w 1 zdaniu (luka G2).
  \item Poprawić \pkCorpusStdCfive{} zdania standardu, które mają 2 przeczenia.
  \item Zalecić legendę przy każdym tekście dla modelu językowego, który zawiera słowo
  \pkq{lub} albo słowo \pkq{albo}.
\end{enumerate}

\paragraph{Dalsze prace.}
Praca zostawia 5 zadań.
\begin{itemize}[leftmargin=1.6em,itemsep=1pt]
  \item Gramatyka formalna dla rdzenia PK.
  \item Sprawdzarka z analizą składni.
  \item Pomiar na modelach o większej liczbie parametrów.
  \item Pomiar na tekstach z kontekstem.
  \item Badanie z udziałem ludzi, które sprawdzi hipotezę o ujawnianiu braków wiedzy.
\end{itemize}
